\documentclass[a4paper,12pt]{article}
\usepackage[T2A]{fontenc}
\usepackage[utf8]{inputenc}
\usepackage[russian]{babel}
\usepackage{amsmath,amssymb,mathtools}
\usepackage{helvet}
% Cyrillic sans-serif: Helvetica lacks T2A glyphs in this environment.
\renewcommand{\sfdefault}{cmss}
\renewcommand{\familydefault}{\sfdefault}
\usepackage{icomma}
\usepackage[left=20mm,right=20mm,top=20mm,bottom=20mm]{geometry}
\usepackage{microtype}
\usepackage{tikz}
\usetikzlibrary{calc,arrows.meta,positioning,shapes.geometric}
\usepackage{tkz-euclide}
\usepackage{enumitem,tabularx,booktabs}
\usepackage{fancyhdr,parskip,setspace}
\usepackage{xcolor,hyperref}
\definecolor{accent}{HTML}{176E72}
\definecolor{darkaccent}{HTML}{18484D}
\definecolor{textgray}{HTML}{303A3C}
\definecolor{softgray}{HTML}{637174}
\definecolor{amber}{HTML}{A86D41}
\hypersetup{colorlinks=true,linkcolor=darkaccent,urlcolor=darkaccent,pdftitle={Чек-лист. Теорема о биссектрисе. 21.10.26},pdfauthor={Лёвин Артём Александрович}}
\setlength{\parindent}{0pt}
\setlength{\parskip}{4pt}
\setstretch{1.17}
\setlist[itemize]{leftmargin=6.5mm,topsep=3pt,itemsep=1pt}
\setlist[enumerate]{leftmargin=8mm,topsep=4pt,itemsep=7pt}
\pagestyle{fancy}
\fancyhf{}
\fancyfoot[C]{\small\color{softgray}\thepage}
\renewcommand{\headrulewidth}{0pt}
\newcommand{\sect}[1]{\par\vspace{2mm}{\Large\bfseries\color{darkaccent}#1}\par\vspace{1mm}{\color{accent!55}\rule{\linewidth}{.5pt}}\par\vspace{1mm}}
\newcommand{\key}[1]{\textbf{\color{darkaccent}#1}}
\newcommand{\ck}[1]{\par\noindent$\square$\enspace #1\par}
\newcommand{\note}[1]{\par\smallskip{\color{darkaccent}\textbf{Проверка.}} #1\par}
\begin{document}
\thispagestyle{empty}
\begin{center}
\vspace*{30mm}
{\small\color{softgray}МАТЕМАТИКА \textbullet{} ОГЭ \textbullet{} ГЕОМЕТРИЯ}\par
\vspace{16mm}
{\fontsize{31}{35}\selectfont\bfseries\color{darkaccent}Чек-лист}\par
\vspace{9mm}
{\LARGE\bfseries Теорема о биссектрисе}\par
\vspace{3mm}
{\large Отношения отрезков и центр вписанной окружности}\par
\vspace{19mm}
{\large София Кальней}\par
\vspace{5mm}
{\normalsize 21 октября 2026 года}\par
\vfill
{\small\color{textgray}Лёвин Артём Александрович эксклюзивно для Софии Кальней}\par
\vspace{17mm}
\end{center}
\begin{tikzpicture}[remember picture,overlay]
\draw[accent!50,line width=1.2pt]
 ([xshift=9mm,yshift=9mm]current page.south west)
 .. controls +(43mm,22mm) and +(-55mm,-6mm) ..
 ([xshift=-9mm,yshift=17mm]current page.south east);
\draw[amber!42,line width=.65pt]
 ([xshift=16mm,yshift=7mm]current page.south west)
 .. controls +(53mm,16mm) and +(-66mm,-7mm) ..
 ([xshift=-16mm,yshift=12mm]current page.south east);
\end{tikzpicture}
\clearpage

\sect{1. Биссектриса делит сторону в особом отношении}
\ck{Помню: биссектриса угла треугольника делит противоположную сторону в отношении двух прилежащих сторон.}
\ck{Правильно сопоставляю отрезки и стороны: верхнему числителю отвечает соответствующая прилежащая сторона.}
\ck{Могу применить теорему и в обратную сторону.}

В $\triangle ABC$ точка $K$ лежит на стороне $BC$, а $AK$ --- биссектриса угла $A$. Тогда
\[
\boxed{\frac{CK}{KB}=\frac{AC}{AB}}.
\]
Обратно: если $K$ лежит \emph{внутри} $BC$ и выполнено это равенство, то $AK$ является биссектрисой $\angle BAC$.

\begin{center}
\begin{tikzpicture}[scale=1.08]
\tkzDefPoint(0,0){A}
\tkzDefPoint(6.7,0){B}
\tkzDefPoint(1.5,3.7){C}
\pgfmathsetmacro{\AC}{sqrt(1.5*1.5+3.7*3.7)}
\pgfmathsetmacro{\q}{\AC/(\AC+6.7)}
\tkzDefPoint({1.5+5.2*\q},{3.7*(1-\q)}){K}
\tkzDrawPolygon[draw=textgray,line width=1pt](A,B,C)
\tkzDrawSegment[draw=accent,line width=1.5pt](A,K)
\tkzDrawPoints[size=2.3,fill=darkaccent](A,B,C,K)
\tkzLabelPoint[below left](A){$A$}
\tkzLabelPoint[below right](B){$B$}
\tkzLabelPoint[above](C){$C$}
\tkzLabelPoint[right](K){$K$}
\tkzMarkAngle[size=.7cm,draw=amber](B,A,K)
\tkzMarkAngle[size=.9cm,draw=amber](K,A,C)
\end{tikzpicture}
\end{center}

\key{Как найти части стороны.} Если $BC=L$, то $CK+KB=L$ и
\[
\boxed{CK=\frac{L\cdot AC}{AC+AB}},\qquad
\boxed{KB=\frac{L\cdot AB}{AC+AB}}.
\]
\key{Пример 1.} $AC=3$, $AB=4$, $BC=5$. Найдём $CK$ и $KB$:
\[
CK=\frac{5\cdot3}{3+4}=\frac{15}{7},\qquad
KB=\frac{5\cdot4}{3+4}=\frac{20}{7}.
\]
\note{$\frac{15}{7}+\frac{20}{7}=5$; $CK:KB=3:4$.}

\key{Пример 2.} $CK=2$, $KB=3$, $AC=4$. По теореме
\[
\frac{2}{3}=\frac{4}{AB}\quad\Longrightarrow\quad\boxed{AB=6}.
\]
\key{Ошибка, которой следует избегать.} Равные части стороны получаются только при равенстве прилежащих сторон. В произвольном треугольнике биссектриса обычно \emph{не совпадает} с медианой или высотой.
\clearpage

\sect{2. Центр вписанной окружности на биссектрисе}
\ck{Знаю, что центр $I$ вписанной окружности --- точка пересечения внутренних биссектрис треугольника.}
\ck{Нахожу треугольник $ACD$, в котором $CI$ и $ID$ --- части стороны $CD$.}
\ck{Использую две пропорции последовательно: сначала в $\triangle ABC$, затем в $\triangle ACD$.}

Пусть $CD$ --- биссектриса $\angle ACB$, $D\in AB$, $I$ --- центр вписанной окружности. Обозначим $AB=c$, $BC=a$, $AC=b$.

\begin{center}
\begin{tikzpicture}[scale=.93]
\tkzDefPoint(0,0){A}
\tkzDefPoint(8,0){B}
\tkzDefPoint(2,5.2){C}
\pgfmathsetmacro{\a}{sqrt(6*6+5.2*5.2)}
\pgfmathsetmacro{\b}{sqrt(2*2+5.2*5.2)}
\pgfmathsetmacro{\per}{\a+\b+8}
\pgfmathsetmacro{\d}{8*\b/(\a+\b)}
\pgfmathsetmacro{\ix}{(8*\b+2*8)/\per}
\pgfmathsetmacro{\iy}{8*5.2/\per}
\tkzDefPoint(\d,0){D}
\tkzDefPoint(\ix,\iy){I}
\tkzDefPoint(\ix,0){H}
\tkzDrawPolygon[draw=textgray,line width=1pt](A,B,C)
\tkzDrawSegments[draw=accent,line width=1.5pt](C,D A,I B,I)
\tkzDrawCircle[draw=amber!80,line width=.9pt](I,H)
\tkzDrawPoints[size=2.4,fill=darkaccent](A,B,C,D,I)
\tkzLabelPoint[below left](A){$A$}
\tkzLabelPoint[below right](B){$B$}
\tkzLabelPoint[above](C){$C$}
\tkzLabelPoint[below](D){$D$}
\tkzLabelPoint[right](I){$I$}
\end{tikzpicture}
\end{center}

\key{Первый шаг:} в $\triangle ABC$ биссектриса $CD$ делит сторону $AB$:
\[
\frac{AD}{DB}=\frac{AC}{BC}=\frac{b}{a},\qquad
AD=\frac{bc}{a+b}.
\]
\key{Второй шаг:} в $\triangle ACD$ луч $AI$ делит угол $CAD$ пополам. Поэтому
\[
\frac{CI}{ID}=\frac{AC}{AD}
=\frac{b}{\frac{bc}{a+b}}.
\]
\key{Итоговая формула:}
\[
\boxed{\frac{CI}{ID}=\frac{a+b}{c}=\frac{AC+BC}{AB}}.
\]
Она верна для любого невырожденного треугольника.

\key{Пример 3.} $AC=7$, $BC=8$, $AB=9$. Тогда
\[
AD=\frac{7\cdot9}{7+8}=\frac{21}{5},\qquad
\frac{CI}{ID}=\frac{7}{21/5}=\frac{35}{21}=\boxed{\frac53}.
\]
\key{Важный вывод.} По неравенству треугольника $AC+BC>AB$. Следовательно,
\[
\boxed{CI>ID.}
\]
Центр вписанной окружности всегда находится ближе к точке $D$, чем к вершине $C$, \emph{если сравнивать две части биссектрисы $CD$}.
\clearpage

\sect{3. Как соединить отношение и периметр}
\ck{Могу перейти от $CI:ID=m:n$ к равенству $AC+BC=(m/n)AB$.}
\ck{Выражаю периметр через $AB$ и нахожу нужную сторону.}
\ck{Умею проверить допустимость данных с помощью неравенств треугольника.}

Если $CI:ID=m:n$, где $m>n>0$, то
\[
\frac{AC+BC}{AB}=\frac mn,\qquad
P=AC+BC+AB=\frac{m+n}{n}\,AB.
\]
Значит,
\[
\boxed{AB=\frac{n}{m+n}\,P},\qquad
\boxed{AC+BC=\frac{m}{m+n}\,P}.
\]
\key{Пример 4.} Центр вписанной окружности делит биссектрису $CD$ в отношении $CI:ID=5:3$. Периметр $P=48$. Найдём $AB$:
\[
AB=\frac{3}{5+3}\cdot48=\boxed{18}.
\]
\note{$AC+BC=48-18=30$ и $30:18=5:3$.}

\sect{4. Типичные ошибки и самопроверка}
\begin{itemize}
\item \key{Перепутать отрезки:} $CI:ID$ --- отношение частей $CD$, а $AD:DB$ --- частей $AB$; у них разные формулы.
\item \key{Использовать одну пропорцию вместо двух:} при поиске $CI:ID$ сначала находят $AD$ через биссектрису $CD$, затем применяют теорему к $AI$ в $\triangle ACD$.
\item \key{Ошибиться с обозначениями:} в общепринятой записи $a=BC$, $b=AC$, $c=AB$. Каждой стороне соответствует противолежащая вершина.
\item \key{Неправильно сократить дробь:} $35/21=5/3$; полезно выполнить проверку общей формулой $(7+8)/9=5/3$.
\item \key{Потерять ограничение:} для настоящего треугольника $AC+BC>AB$, поэтому $CI:ID$ всегда больше $1$.
\end{itemize}
\clearpage

\begin{center}
{\Large\bfseries\color{darkaccent}Алгоритм: найти $CI:ID$}\par
\vspace{3mm}
{\small\color{softgray}Треугольник $ABC$; $CD$ --- биссектриса; $I$ --- центр вписанной окружности}
\end{center}
\vspace{14mm}
\begin{center}
\begin{tikzpicture}[
 >=Stealth,node distance=10mm,
 start/.style={ellipse,draw=darkaccent,line width=.85pt,minimum height=12mm,align=center,text width=9.1cm,fill=accent!4},
 flowstep/.style={rectangle,rounded corners=1mm,draw=accent,line width=.9pt,align=center,text width=11.2cm,minimum height=16mm,inner sep=3mm},
 arr/.style={->,draw=darkaccent,line width=1.1pt}
]
\node[start](a){Дано: $\triangle ABC$, биссектриса $CD$, центр $I$};
\node[flowstep,below=of a](b){Проведите вторую биссектрису $AI$; выделите два треугольника: $ABC$ и $ACD$};
\node[flowstep,below=of b](c){В $\triangle ABC$ запишите $\displaystyle\frac{AD}{DB}=\frac{AC}{BC}$};
\node[flowstep,below=of c](d){Учтите $AD+DB=AB$ и выразите $\displaystyle AD=\frac{AB\cdot AC}{AC+BC}$};
\node[flowstep,below=of d](e){В $\triangle ACD$ используйте биссектрису $AI$:\\ $\displaystyle\frac{CI}{ID}=\frac{AC}{AD}$};
\node[flowstep,below=of e](f){Подставьте $AD$ и сократите:\\ $\displaystyle\frac{CI}{ID}=\frac{AC+BC}{AB}$};
\node[start,below=of f](g){Проверьте, что $CI:ID>1$};
\draw[arr](a)--(b);
\draw[arr](b)--(c);
\draw[arr](c)--(d);
\draw[arr](d)--(e);
\draw[arr](e)--(f);
\draw[arr](f)--(g);
\end{tikzpicture}
\end{center}
\clearpage

\sect{5. Самостоятельная тренировка --- Путь А}
\textbf{10 заданий.} Все треугольники невырожденные. В задачах 1--3 точка $K\in BC$, а $AK$ --- биссектриса $\angle A$. В задачах 4--10 точка $D\in AB$, $CD$ --- биссектриса $\angle C$, $I$ --- центр вписанной окружности. Ответы приведены на следующей странице.
\begin{enumerate}
\item Дано $AC=6$, $AB=9$. Найдите $CK:KB$.
\item Дано $AC=6$, $AB=10$, $BC=8$. Найдите $CK$ и $KB$.
\item Дано $CK=3$, $KB=5$, $AC=9$. Найдите $AB$.
\item Дано $AC=7$, $BC=8$, $AB=9$. Найдите $CI:ID$.
\item Дано $AC=10$, $BC=13$, $AB=15$. Найдите $CI/ID$ и объясните, почему отношение больше единицы.
\item Известно $CI:ID=7:4$ и $AB=12$. Найдите $AC+BC$ и периметр треугольника.
\item Известно $CI:ID=5:3$ и периметр треугольника равен $64$. Найдите сторону $AB$.
\item Дано $AB=18$, $AC=15$, $BC=12$. Найдите $AD$, $DB$ и $CI:ID$.
\item Докажите, что в любом невырожденном треугольнике $CI>ID$. Может ли выполняться $CI=ID$?
\item Дано $AB=24$, $AD=9$, $AC=12$. Найдите $BC$ и $CI:ID$.
\end{enumerate}
\clearpage

\sect{6. Ответы к самостоятельной тренировке}
\renewcommand{\arraystretch}{1.45}
\begin{tabularx}{\linewidth}{@{}c>{\raggedright\arraybackslash}p{44mm}X@{}}
\toprule
\textbf{№} & \textbf{Ответ} & \textbf{Ключевой шаг / проверка}\\
\midrule
1 & $2:3$ & $CK:KB=AC:AB=6:9$.\\
2 & $CK=3$, $KB=5$ & $CK=8\cdot6/(6+10)$; сумма равна $8$.\\
3 & $AB=15$ & $3/5=9/AB$.\\
4 & $5:3$ & $(AC+BC)/AB=15/9=5/3$.\\
5 & $\dfrac{23}{15}$ & $10+13=23>15$.\\
6 & $AC+BC=21$, $P=33$ & $AC+BC=12\cdot7/4$.\\
7 & $AB=24$ & $AB=64\cdot3/(5+3)$.\\
8 & $AD=10$, $DB=8$, $CI:ID=3:2$ & $AD=18\cdot15/27$; $CI/ID=27/18$.\\
9 & $CI>ID$ всегда; равенство невозможно & $\dfrac{CI}{ID}=\dfrac{AC+BC}{AB}>1$ по неравенству треугольника.\\
10 & $BC=20$, $CI:ID=4:3$ & $DB=24-9=15$; $BC=15\cdot12/9=20$.\\
\bottomrule
\end{tabularx}

\vspace{8mm}
\key{Контрольный ориентир:} в задачах о центре вписанной окружности сравнивайте найденное отношение с единицей. Ответ $CI:ID\le 1$ указывает на ошибку вычислений или некорректные исходные данные.
\end{document}
